\documentclass[11pt]{article}
\usepackage[T1]{fontenc}
\usepackage{lmodern,amsmath,amssymb,amsthm,mathtools}
\usepackage[margin=1in]{geometry}
\usepackage[colorlinks=true,linkcolor=blue,citecolor=blue,urlcolor=blue]{hyperref}
\hypersetup{pdftitle={An explicit complex polynomial retract in five variables: a consequence of affine-space cancellation failure},pdfauthor={Alec Kriebel}}
\newtheorem{theorem}{Theorem}
\newtheorem{proposition}[theorem]{Proposition}
\newtheorem{lemma}[theorem]{Lemma}
\newcommand{\C}{\mathbb C}
\newcommand{\A}{\mathbb A}
\newcommand{\ev}{\operatorname{ev}}
\newcommand{\Spec}{\operatorname{Spec}}
\title{An explicit complex polynomial retract in five variables:\\a consequence of affine-space cancellation failure}
\author{Alec Kriebel\\\small ORCID: \href{https://orcid.org/0009-0001-9320-500X}{0009-0001-9320-500X}}
\date{October 6, 2026}
\begin{document}
\maketitle
\begin{abstract}
OpenAI's explicit complex cancellation fourfold has a polynomial
one-variable cylinder and is not affine four-space. We transport coefficient
inclusion and evaluation at zero through its concrete cylinder coordinates.
This gives a complex algebra retract of a five-variable polynomial ring,
of transcendence degree four, that is not a polynomial algebra. We provide
split maps and a five-component idempotent polynomial endomorphism, in
factored form with no denominators, and exact symbolic certificates for
the coordinate identities, including the fiber at the parameter zero.
The implication from cancellation failure to a negative answer to Costa's
retract question is classical; the nonpolynomiality input and cancellation
breakthrough belong to OpenAI. This note makes that consequence explicit
in ambient dimension five and makes no assertion about ambient dimension four.
\end{abstract}

\section{The question and the inherited input}
Costa asked whether a retract of $k^{[n]}$ containing the field $k$ must
itself be a polynomial $k$-algebra \cite{Costa}. Here a retract is specified
by unital $k$-algebra maps $i:A\to B$ and $r:B\to A$ with $r i=1_A$;
we identify $A$ with $i(A)\subset B$. Nagamine explicitly records the
classical implication: a positive answer in $n+1$ variables implies
cancellation for $\A^n$ \cite[Proposition 1.4]{Nagamine}. His proof is
coefficient inclusion followed by evaluation of the cylinder variable at zero.
Nagamine settles the ambient three-variable case in characteristic zero.
Earlier literature also states the polynomial-extension retract relationship
\cite[Introduction]{EN}. More generally, every characteristic-zero polynomial-ring retract of
transcendence degree two is polynomial \cite[Theorem 5.8]{CDDG}.
The August 2026 version of Chakraborty--Pal still records the general
characteristic-zero question as open for ambient $n\geq4$ \cite{CP}.

We use the following particular input, with its authorship explicit.
Set
\begin{equation}\label{eq:A}
\begin{gathered}
P=\C[p,s,u,F,J],\quad x=s^2+u^3+p^2F,\\
H=x^2F-(1+2sx)J-p^2J^2-pu,\qquad A=P/(H).
\end{gathered}
\end{equation}
OpenAI \cite[Theorem 1 and Sections 2--6]{OAI} proves that $A$ is an
integral four-dimensional complex algebra, that $A[w]\cong\C^{[5]}$,
and that $A\not\cong\C^{[4]}$. The pinned source commit is
\begin{quote}\small\ttfamily adc7f1241b42e322a6451854ab7e4b4c146bf78a\end{quote}
The source manuscript is
dated September 23, 2026; the public repository evidence checked for this
note dates from October 6, 2026. We do not equate the manuscript date with
a verified public disclosure date.

The written nonpolynomiality proof is a pivotal external theorem, rather
than a conclusion of the computations below. Its mechanism is worth
stating precisely. The valuation degeneration is the smooth algebra
\[
G=R[\tau,I/\tau^2],\quad
R=\C[x,y,z,u]/(xy-z(z+1)),
\]
where $s_R=y-x(x-u^3)$ and
$I=(x-s_R^2-u^3,z-s_Rx)$. Polynomiality would give a nonzero homogeneous
locally nilpotent derivation of $G$ with a positive-degree invariant.
The determinant torus bundle
\[
\widetilde R=\C[a,d,b,c,u]/(ac-bd-1),\quad
(x,y,z)=(ab,dc,db)
\]
principalizes $I$ to $(v)$, where $v=a^3b-a^2u^3-d^2$. A normalized
additive lift to its line-bundle pullback, together with the invariant,
produces a nonzero locally nilpotent derivation $E$ of $\widetilde R$
with nonpositive shift for weights $(1,1,-1,-1,0)$ and $E^2(v)=0$.
The highest part for auxiliary weights $(0,3,-3,0,2)$ satisfies
$E'^2(d^2+a^2u^3)=0$. The Mason--Stothers orbit-polynomial argument
and the nonsquareness of $-u^3$ in $\C(a,d,b)(u)$ force $E'$ to fix
$a,d,u$. Regularity across the determinant charts then forces its
nonzero shift to be positive, a contradiction. The full proofs of the
exact degeneration, bundle lift, sign extraction, and rigidity are in
\cite[Sections 3--6]{OAI}; our dependency audit checks each of these steps.

\begin{theorem}\label{thm:main}
The algebra $A$ in \eqref{eq:A} is a smooth integral complex algebra of
transcendence degree four. There are explicit complex algebra maps
$i:A\to B=\C[p,s,u,m,e]$ and $r:B\to A$ with $r i=1_A$.
The image $i(A)$ is not a polynomial complex algebra. The five polynomial
components given in Section~\ref{sec:idempotent} define the idempotent
algebra endomorphism $i r$ of $B$ and hence an idempotent endomorphism
of $\A^5$ with image isomorphic to $\Spec A$.
\end{theorem}
Thus Costa's assertion is false over $\C$ in ambient dimension five.
The four-variable retract question is not settled here. The new work
in this note is the explicit transport, formulas, and verification
package; the conditional reduction and the cancellation input are inherited.

\section{Concrete cylinder coordinates}
All substitutions in this section are polynomial algebra homomorphisms.
We reproduce the coordinate construction of \cite[Section 2]{OAI} in
a form adapted to the split maps.

For independent $p,s,u,F,J,w$, with $x$ as in \eqref{eq:A}, define
\begin{equation}\label{eq:phi}
\begin{gathered}
D_w=-3u^2w+3p^2uw^2-p^4w^3,\\
\Phi(p)=p,\quad \Phi(w)=w,\quad
\Phi(u)=u-p^2w,\quad \Phi(s)=s+p^2xD_w,\\
\Phi(F)=F-(1+2sx)D_w-p^2x^2D_w^2,\quad
\Phi(J)=J-x^2D_w.
\end{gathered}
\end{equation}
The polynomial $x$ is fixed. Replacing $w$ by $-w$ in these substitution
formulas gives the inverse. One way to verify them is to exponentiate
the locally nilpotent derivation
\[
\Delta(p)=0,\quad\Delta(u)=-p^2,\quad
\Delta(s)=-3p^2xu^2,\quad
\Delta(F)=(6sx+3)u^2,\quad\Delta(J)=3x^2u^2.
\]
Indeed, in the localized coordinates
$y=s+x(x-u^3)$, $z=sx+p^2J$, the derivation is
$-p^2\partial_u$, fixing $p,x,y,z$. Its displayed values are in $P$,
so it restricts to a locally nilpotent derivation of $P$.
Direct substitution gives $\Phi(H)=H+p^3w$. Consequently
\begin{equation}\label{eq:T}
\overline\Phi:A[w]\longrightarrow
T=P[w]/(H+p^3w)
\end{equation}
is an isomorphism, with inverse induced by $\Phi^{-1}$.

Put $x_0=s^2+u^3$ and introduce
\begin{equation}\label{eq:LM}
L=x_0^2F-(1+2sx_0)J,\qquad
M=(1-2sx_0)F+4s^2J.
\end{equation}
The determinant of this linear coordinate change in $F,J$ is one;
its inverse is
\[
f(L)=4s^2L+(1+2sx_0)M,\qquad
j(L)=(2sx_0-1)L+x_0^2M.
\]
In the coefficient ring $K_0=\C[p,s,u,M]$, set
\begin{equation}\label{eq:Q}
\begin{gathered}
Q(L)=2x_0f(L)^2-2sf(L)j(L)-j(L)^2,\\
C=Q(0)=M^2(2x_0+6sx_0^2+4s^2x_0^3-x_0^4),\\
Q(L)=C+LQ_1(L),\qquad L_*=pu-p^2C.
\end{gathered}
\end{equation}
The factor $Q_1$ is a polynomial. An explicit expression, avoiding any
division in the formulas, is $Q_1(L)=q_1+q_2L$, where
\begin{equation}\label{eq:qcoeff}
\begin{gathered}
a_0=4s^2,\quad b_0=1+2sx_0,\quad c_0=2sx_0-1,\quad d_0=x_0^2,\\
q_1=M\{4x_0a_0b_0-2s(a_0d_0+b_0c_0)-2c_0d_0\},\\
q_2=2x_0a_0^2-2sa_0c_0-c_0^2.
\end{gathered}
\end{equation}
In these coordinates the equation is
\begin{equation}\label{eq:Hexp}
H(L)=L-pu+p^2Q(L)+p^4f(L)^3.
\end{equation}
Define the following polynomial in $K_0[L,w]$:
\begin{equation}\label{eq:e0}
h=u-pQ(L)-p^3f(L)^3-p^2w,\qquad
e_0=-w-pf(L)^3-Q_1(L)h.
\end{equation}
Two exact identities are
\begin{align}
H(L_*)&=p^3\{(u-pC)Q_1(L_*)+pf(L_*)^3\},\label{eq:root}\\
p^3e_0-(L-L_*)&=(p^2Q_1(L)-1)(H(L)+p^3w).\label{eq:cert}
\end{align}
In particular the class of $e_0$ in $T$ is the unlocalized polynomial
representative of $(L-L_*)/p^3$; no division is required at $p=0$.

Define $B=\C[p,s,u,m,e]$, replacing $M$ by $m$ in expressions evaluated
in $B$, and put
\begin{equation}\label{eq:LW}
\begin{gathered}
\ell=pu-p^2C+p^3e,\quad f=f(\ell),\quad j=j(\ell),\\
W=-e-(u-pC+p^2e)Q_1(\ell)-pf^3.
\end{gathered}
\end{equation}
These satisfy $H(\ell)+p^3W=0$ identically in $B$.
There are maps
\[
q:B\longrightarrow T,\quad (p,s,u,m,e)\longmapsto(p,s,u,M,e_0),
\]
and $g:T\to B$ with $(p,s,u,F,J,w)\mapsto(p,s,u,f,j,W)$.

\begin{lemma}\label{lem:inverse}
The maps $q$ and $g$ are inverse isomorphisms.
\end{lemma}
\begin{proof}
The determinant-one calculation recovers $L,M$ from $f(L),j(L)$.
Equations \eqref{eq:LW} and \eqref{eq:cert} recover $L$ in $T$.
Substituting $(L,w)=(\ell,W)$ into \eqref{eq:e0} gives exactly $e$,
as a polynomial identity, and hence $gq=1_B$ on all five generators.
For the other composition, $qg$ fixes $p,s,u,M$ and $L$.
The defining equation and $H(L_*)\in p^3K_0$ then give
$p^3(qg(w)-w)=0$. The ring $T$ is a domain and $p\ne0$:
$T\cong A[w]$, and the UFD localization argument below proves that $A$
is a domain with nonzero $p$. Cancellation proves $qg(w)=w$.
Thus $qg$ fixes every generator of $T$.
\end{proof}

For clarity, the domain argument does not use nonpolynomiality.
Over $\C[p,p^{-1}]$, $x,y,z,u$ are coordinates, with inverse
\[
s=y-x(x-u^3),\quad F=p^{-2}(x-s^2-u^3),\quad J=p^{-2}(z-sx).
\]
Here $H=p^{-2}(xy-z(z+1))-pu$ is a coordinate, with a unit coefficient
of $u$, so $(H)$ is prime after localization. Since $p$ does not divide
$H$ (its residue is the nonzero polynomial $L$), no factor of $H$ can
be lost by inverting $p$. Unique factorization proves $H$ irreducible
and prime in $P$. Moreover $A[p^{-1}]=\C[p,p^{-1},x,y,z]$,
so $A$ has transcendence degree four.

\section{Split maps and the five components}\label{sec:idempotent}
Let $\jmath:A\to A[w]$ be coefficient inclusion and
$\epsilon:A[w]\to A$ send $w$ to zero. The actual cylinder isomorphism
is $\theta=g\overline\Phi$. Set
\begin{equation}\label{eq:split}
i=\theta\jmath,\qquad r=\epsilon\overline\Phi^{-1}q.
\end{equation}
The retraction is particularly simple on generators:
\begin{equation}\label{eq:r}
r(p)=p,\quad r(s)=s,\quad r(u)=u,\quad r(m)=M,\quad
r(e)=E_A:=-pF^3-Q_1(L)(u-pQ(L)-p^3F^3)
\end{equation}
in $A$, where $L,M,Q,Q_1$ are evaluated at the original $s,u,F,J$.
This follows because every exponential correction vanishes at $w=0$.

For the inclusion use \eqref{eq:LW}, now in $B$, and set
\begin{equation}\label{eq:inclusion}
\begin{gathered}
X=s^2+u^3+p^2f,\quad D=-3u^2W+3p^2uW^2-p^4W^3,\\
S=s+p^2XD,\quad U=u-p^2W,\\
\mathcal F=f-(1+2sX)D-p^2X^2D^2,\quad
\mathcal J=j-X^2D.
\end{gathered}
\end{equation}
Then $i(p,s,u,F,J)=(p,S,U,\mathcal F,\mathcal J)$.
In particular $H(p,S,U,\mathcal F,\mathcal J)=0$.
For arbitrary arguments define the two polynomial functions
\[
\mathcal M(s,u,F,J)=(1-2s(s^2+u^3))F+4s^2J,
\]
and $\mathcal E(p,s,u,F,J)=E_A$ from \eqref{eq:r}, using
\eqref{eq:LM}--\eqref{eq:qcoeff}. The promised five components are
\begin{equation}\label{eq:rho}
\boxed{\rho=(p,S,U,
\mathcal M(S,U,\mathcal F,\mathcal J),
\mathcal E(p,S,U,\mathcal F,\mathcal J)).}
\end{equation}
This is a fully specified polynomial substitution on $B$, with all
abbreviations defined above. It represents $i r$.

\begin{proof}[Proof of Theorem~\ref{thm:main}]
Lemma~\ref{lem:inverse} and \eqref{eq:T} prove that $\theta$ is an
actual isomorphism. Hence
$r i=\epsilon\theta^{-1}\theta\jmath=\epsilon\jmath=1_A$ and
$(i r)^2=i(r i)r=i r$. Its image is exactly $i(A)$, since $r$ is
surjective and $i$ is injective. The cited nonpolynomiality theorem
\cite[Theorem 1]{OAI} gives $A\not\cong\C^{[4]}$; since $A$ has
transcendence degree four, it is not a polynomial algebra in any number
of variables.

Smoothness can be checked directly. On $p\ne0$ the displayed Laurent
polynomial presentation is smooth. On $p=0$, the partial derivatives
of $H$ with respect to $F,J$ are $x_0^2$ and $-(1+2sx_0)$.
They cannot vanish simultaneously, since
$4s^2x_0^2+(1-2sx_0)(1+2sx_0)=1$. The hypersurface Jacobian criterion
therefore proves smoothness everywhere.

On spectra, $\pi=\Spec(i):\A^5\to\Spec A$ and
$\sigma=\Spec(r):\Spec A\to\A^5$ obey
$\pi\sigma=1_{\Spec A}$. The endomorphism with components \eqref{eq:rho}
is $\sigma\pi$. Its scheme-theoretic image is the closed subscheme
$\sigma(\Spec A)$, with coordinate algebra $B/\ker r\cong A$.
This explains both the contravariance and the image assertion.
\end{proof}

The zero fiber is also explicit. Let $\kappa=q_1/m$, meaning the
bracketed polynomial in \eqref{eq:qcoeff}, not a division operation.
At $p=0$, $W=-e-u\kappa m$ and
\[
\rho(0,s,u,m,e)=(0,s,u,m+3u^2W,-u\kappa(m+3u^2W)).
\]
Its new $W$ is zero, so a second application fixes all components.
This verifies the limiting fiber independently of localization.

\section{Verification, attribution, and limits}
The supplement identifies all source hashes, the exact dependency
statements, audit scopes, and reproducible checks. The symbolic script
uses integer polynomial arithmetic through SymPy. It checks the
derivation, exponential inverse, determinant-one change, root and
fifth-coordinate certificates, and the polynomial inverse identities.
The other inverse composition on $w$ uses the proved domain and
$p$-cancellation, rather than a sample-point test. The structural proof
of \eqref{eq:split} then certifies idempotence and the image algebra.
Additional exact rational examples are regression checks only; they do
not prove nonpolynomiality or replace polynomial identities.

The actual upstream Lean declaration is
\texttt{OAI.ComplexCancellation.main} in
\path{OAI/Algebra/AffineCancellation/Main.lean}; its statement matches
\eqref{eq:A} and the cylinder/nonpolynomiality claim. The comparator's
\texttt{sorry} is a separate challenge statement. We inspected the
actual solution sources and their semantic scope, but did not reproduce
the full Lean kernel build: the local exact-version dependency build
was stopped by disk limits. Accordingly this note relies on the
independently audited written theorem, and claims neither a reproduced
formal proof nor a formalization of the follow-on retract theorem.

Our primary-source and companion-corpus audit found no earlier explicit
statement or five-component transport for this particular complex
example. A failed search is not proof of priority. The general
cylinder-to-retract mechanism is already explicit in Nagamine and is
not claimed as new. We attribute the defining example, stabilization,
and nonpolynomiality to OpenAI using its manuscript-specific citation.
This is a consequence note, not an independent cancellation breakthrough.

\paragraph{Research disclosure.}
AI tools were used extensively for research, drafting, symbolic
verification, and adversarial audits. The upstream work is itself
AI-produced and is credited to its supplied author, OpenAI.
This preprint has not undergone conventional human peer review or
refereeing at publication. Automated adversarial reviews are not human
peer review. No external individuals were contacted.

\begin{thebibliography}{9}
\bibitem{Costa} D.~L.~Costa, \emph{Retracts of polynomial rings},
J. Algebra \textbf{44} (1977), 492--502.
\href{https://doi.org/10.1016/0021-8693(77)90197-1}{doi:10.1016/0021-8693(77)90197-1}.
\bibitem{Nagamine} T.~Nagamine, \emph{A note on retracts of polynomial
rings in three variables}, J. Algebra \textbf{534} (2019), 339--343;
\href{https://arxiv.org/abs/1811.04153v2}{arXiv:1811.04153v2} (2018).
\bibitem{CDDG} S.~Chakraborty, N.~Dasgupta, A.~K.~Dutta and N.~Gupta,
\emph{Some results on retracts of polynomial rings},
\href{https://arxiv.org/abs/1910.11023v1}{arXiv:1910.11023v1} (2019).
\bibitem{CP} S.~Chakraborty and M.~Pal, \emph{Monomial and binomial
retracts of polynomial rings},
\href{https://arxiv.org/abs/2504.14382v2}{arXiv:2504.14382v2} (2026).
\bibitem{EN} N.~Epstein and H.~D.~Nguyen, \emph{Algebra retracts and
Stanley--Reisner rings}, J. Pure Appl. Algebra \textbf{218} (2014),
1665--1682; \href{https://arxiv.org/abs/1301.3967v1}{arXiv:1301.3967v1} (2013).
\bibitem{OAI} OpenAI, \emph{An explicit failure of complex affine-space
cancellation}, OpenAI Math Release preprint, 2026.
\href{https://github.com/openai/math/blob/adc7f1241b42e322a6451854ab7e4b4c146bf78a/preprints/An-explicit-failure-of-complex-affine-space-cancellation-September-23-2026/paper.pdf}{Pinned manuscript and repository}.
\end{thebibliography}
\end{document}
